% corpus_supplement: Erdős problem #458
% source: https://www.erdosproblems.com/latex/458
% fetched_at: 2026-07-14T18:48:17Z
% payload_sha256: 87fce800f175bd309d36f399813126930c1e5deac216b370ff68f4d537bce3a1
% statement/remarks/references extracted by erdos_ingest.extract_source_page;
% rendered by erdos_ingest.render_tex — do not edit
\documentclass{article}
\usepackage{amsmath, amssymb, amsthm}
\usepackage{hyperref}
\usepackage{geometry}
\geometry{margin=1in}

\title{Erd\H{o}s Problem \#458}
\date{Last updated: 2025-08-31}
\author{Source: erdosproblems.com}

\begin{document}
\maketitle

\noindent\textbf{Status:} OPEN \quad \textbf{No prize}

\noindent\textbf{Tags:} number theory, primes

\medskip
\noindent\textbf{Problem Statement:}

Let $[1,\ldots,n]$ denote the least common multiple of $\{1,\ldots,n\}$. Is it true that, for all $k\geq 1$,\[[1,\ldots,p_{k+1}-1]< p_k[1,\ldots,p_k]?\]

\bigskip
\noindent\textbf{Remarks:}

Erd\H{o}s and Graham write this is 'almost certainly' true, but the proof is beyond our ability, for (at least) two reasons:
{UL}
{LI}Firstly, one has to rule out the possibility of many primes $q$ such that $p_k<q^2<p_{k+1}$. There should be at most one such $q$, which would follow from $p_{k+1}-p_k<p_k^{1/2}$, which is essentially the notorious Legendre's conjecture.{/LI}
{LI} The small primes also cause trouble.{/LI}
{/UL}

\bigskip
\noindent\small{Source: \url{https://www.erdosproblems.com/458}}
\end{document}
